\begin{abstract}
At the nominal false-alarm level $\alpha=0.20$, a target-only leave-one-image-out (LOIO) alarm reaches an empirical false-alarm rate (FAR) of 0.341 on Gaussian-corrupted MVTec at $k=4$, approximately 1.7 times the requested level. This motivates a study-design question: how much source evidence, measured in images or category units, can support threshold certification for a new category? Using a frozen DINOv2 principal component analysis (PCA) residual ranker, we audit 15 MVTec and 12 VisA categories under four corruptions. Target-only rank values have the finite resolution floor $1/(k{+}1)$, yet crossing it does not restore false-alarm stability under shift. We derive a category-count feasibility calculus. Even with zero observed category losses, an optimistic multiplicity-free, deterministic, uniformly valid, distribution-free 95\% upper confidence bound (UCB) requires at least 14, 29, and 59 independent and identically distributed (iid) category draws at $\alpha=0.20$, 0.10, and 0.05. These counts are necessary, not sufficient. Cross-category Reliability Estimation with Source Support (CRESS) separates reference, threshold-proposal, and certification categories. With only three or four certification categories, every target cell evaluated under category-unit certification across 960 frozen gate configurations receives the fail-closed threshold $\tau^\star=0$; the smallest category-level UCB is 0.950, far above the largest tested $\alpha$. Under an idealized iid source-image model, image-unit sensitivity bounds select positive thresholds in 36.7\% to 60.3\% of target cells, but address the selected source mixture rather than marginal risk over a new-category draw. The main contribution is a quantitative feasibility boundary and an estimand-aware protocol for deciding when source evidence can, and cannot, justify a transferable reliability claim.
\end{abstract}
